%%%%%%%%%%%%%%%%%%%%%%%%%%%%%%%%%%%%%%%%%%%%%%%%%%%%%%%%%%%%%%%%%%%%%%%%%%%%%%%%
%%%%%%%%%%%%%%%%%%%%%%%%%%%%%%%%%%%%%%%%%%%%%%%%%%%%%%%%%%%%%%%%%%%%%%%%%%%%%%%%
%%%%%%%%%%%%%%%%%%%%%%%%%%%%%%%%%%%%%%%%%%%%%%%%%%%%%%%%%%%%%%%%%%%%%%%%%%%%%%%%
%%%%%%%%%%%%%%%%%%%%%%%%%%%%%%%%%%%%%%%%%%%%%%%%%%%%%%%%%%%%%%%%%%%%%%%%%%%%%%%%
\chapter{Décomposition en éléments simples\label{annexe-DES}}
%%%%%%%%%%%%%%%%%%%%%%%%%%%%%%%%%%%%%%%%%%%%%%%%%%%%%%%%%%%%%%%%%%%%%%%%%%%%%%%%
%%%%%%%%%%%%%%%%%%%%%%%%%%%%%%%%%%%%%%%%%%%%%%%%%%%%%%%%%%%%%%%%%%%%%%%%%%%%%%%%
%%%%%%%%%%%%%%%%%%%%%%%%%%%%%%%%%%%%%%%%%%%%%%%%%%%%%%%%%%%%%%%%%%%%%%%%%%%%%%%%
%%%%%%%%%%%%%%%%%%%%%%%%%%%%%%%%%%%%%%%%%%%%%%%%%%%%%%%%%%%%%%%%%%%%%%%%%%%%%%%%

%%%%%%%%%%%%%%%%%%%%%%%%%%%%%%%%%%%%%%%%%%%%%%%%%%%%%%%%%%%%%%%%%%%%%%%%%%%%%%%%
%%%%%%%%%%%%%%%%%%%%%%%%%%%%%%%%%%%%%%%%%%%%%%%%%%%%%%%%%%%%%%%%%%%%%%%%%%%%%%%%
%%%%%%%%%%%%%%%%%%%%%%%%%%%%%%%%%%%%%%%%%%%%%%%%%%%%%%%%%%%%%%%%%%%%%%%%%%%%%%%%
\section{Contexte}
%%%%%%%%%%%%%%%%%%%%%%%%%%%%%%%%%%%%%%%%%%%%%%%%%%%%%%%%%%%%%%%%%%%%%%%%%%%%%%%%
%%%%%%%%%%%%%%%%%%%%%%%%%%%%%%%%%%%%%%%%%%%%%%%%%%%%%%%%%%%%%%%%%%%%%%%%%%%%%%%%
%%%%%%%%%%%%%%%%%%%%%%%%%%%%%%%%%%%%%%%%%%%%%%%%%%%%%%%%%%%%%%%%%%%%%%%%%%%%%%%%
En automatique, la détermination d'une réponse temporelle $s(t)$
correspond à déterminer la transformée 
de Laplace inverse d'une fraction rationnelle $S(p)$ 
définie dans le domaine de Laplace. Autrement dit, 
\[
s(t)=\laplacei{S(p)}
\]
Cette inversion passe généralement par l'utilisation des tables de 
transformées de Laplace (cf. \Cref{annexe-lap}).
Lorsque ces tables ne sont pas complètes, on utilise alors 
la~\gls{des} de $S(p)$ permettant de réécrire cette fraction 
rationnelle en une somme de fractions rationnelles 
usuellement présentes dans les tables.

Dans cette annexe, nous présenterons les techniques de~\gls{des} les plus 
couramment rencontrées dans l'étude des~\gls{slci}. 
Cette présentation n'est pas exhaustive, et ne remplacera donc 
pas la lecture du chapitre du cours de mathématiques qui lui est consacré.

%%%%%%%%%%%%%%%%%%%%%%%%%%%%%%%%%%%%%%%%%%%%%%%%%%%%%%%%%%%%%%%%%%%%%%%%%%%%%%%%
%%%%%%%%%%%%%%%%%%%%%%%%%%%%%%%%%%%%%%%%%%%%%%%%%%%%%%%%%%%%%%%%%%%%%%%%%%%%%%%%
%%%%%%%%%%%%%%%%%%%%%%%%%%%%%%%%%%%%%%%%%%%%%%%%%%%%%%%%%%%%%%%%%%%%%%%%%%%%%%%%
\section{Fractions rationnelles rencontrées en automatique}
%%%%%%%%%%%%%%%%%%%%%%%%%%%%%%%%%%%%%%%%%%%%%%%%%%%%%%%%%%%%%%%%%%%%%%%%%%%%%%%%
%%%%%%%%%%%%%%%%%%%%%%%%%%%%%%%%%%%%%%%%%%%%%%%%%%%%%%%%%%%%%%%%%%%%%%%%%%%%%%%%
%%%%%%%%%%%%%%%%%%%%%%%%%%%%%%%%%%%%%%%%%%%%%%%%%%%%%%%%%%%%%%%%%%%%%%%%%%%%%%%%
Dans le cas qui nous intéresse la fraction rationnelle est la 
réponse $S(p)$ définie dans le domaine de Laplace.
Cette grandeur est de la forme,
\[
S(p)=\dfrac{N(p)}{D(p)}
\]
où $N(p)$ et $D(p)$ sont deux polynômes de degrés $m$ et $n$ respectivement.
En général, nous aurons à faire à des systèmes pour lesquels $n>m$. 
Une des conséquences est que \textbf{la décomposition en éléments 
simples ne comportera pas de partie entière.}

%%%%%%%%%%%%%%%%%%%%%%%%%%%%%%%%%%%%%%%%%%%%%%%%%%%%%%%%%%%%%%%%%%%%%%%%%%%%%%%%
%%%%%%%%%%%%%%%%%%%%%%%%%%%%%%%%%%%%%%%%%%%%%%%%%%%%%%%%%%%%%%%%%%%%%%%%%%%%%%%%
%%%%%%%%%%%%%%%%%%%%%%%%%%%%%%%%%%%%%%%%%%%%%%%%%%%%%%%%%%%%%%%%%%%%%%%%%%%%%%%%
\section{Division euclidienne : cas $\deg N\ge\deg D$}
%%%%%%%%%%%%%%%%%%%%%%%%%%%%%%%%%%%%%%%%%%%%%%%%%%%%%%%%%%%%%%%%%%%%%%%%%%%%%%%%
%%%%%%%%%%%%%%%%%%%%%%%%%%%%%%%%%%%%%%%%%%%%%%%%%%%%%%%%%%%%%%%%%%%%%%%%%%%%%%%%
Si $n\le m$, il faut procéder à une division euclidienne de la 
fraction rationnelle telle que:
\[
  N(p) = Q(p)D(p) + R(p) 
\]
Avec $Q(p)$ et $R(p)$ respectivement les polynômes quotient et reste et tel que
$\text{deg}R(p)<\text{deg}D(p)$ ou encore :
\[
    R(p) = N(p) - Q(p)D(p)
\]

Prenons l'exemple d'une fraction rationnelle donnée par:
\[
    \dfrac{p(p-1)}{(p+1)^2}
\]
Pour déterminer $Q(p)$, il suffit de prendre le rapport du terme de plus haut degré 
de $N$ par celui de $D$, ici $Q$ étant égal à 1:
\[
    R(p)=p(p-1)-(p+1)^2=-3p-1
\]
Donc:
\[
    N(p)=(p+1)^2-3(p+1)+2
\]
et 
\[
    S(p)=1-\dfrac{3}{(p+1)}+\dfrac{2}{(p+1)^2}
\]
%%%%%%%%%%%%%%%%%%%%%%%%%%%%%%%%%%%%%%%%%%%%%%%%%%%%%%%%%%%%%%%%%%%%%%%%%%%%%%%%
%%%%%%%%%%%%%%%%%%%%%%%%%%%%%%%%%%%%%%%%%%%%%%%%%%%%%%%%%%%%%%%%%%%%%%%%%%%%%%%%
%%%%%%%%%%%%%%%%%%%%%%%%%%%%%%%%%%%%%%%%%%%%%%%%%%%%%%%%%%%%%%%%%%%%%%%%%%%%%%%%
\section{Décomposition en éléments simples}
%%%%%%%%%%%%%%%%%%%%%%%%%%%%%%%%%%%%%%%%%%%%%%%%%%%%%%%%%%%%%%%%%%%%%%%%%%%%%%%%
%%%%%%%%%%%%%%%%%%%%%%%%%%%%%%%%%%%%%%%%%%%%%%%%%%%%%%%%%%%%%%%%%%%%%%%%%%%%%%%%
%%%%%%%%%%%%%%%%%%%%%%%%%%%%%%%%%%%%%%%%%%%%%%%%%%%%%%%%%%%%%%%%%%%%%%%%%%%%%%%%
Soit $S(p)=\dfrac{N(p)}{D(p)}$ une fraction rationnelle. On considère la 
décomposition de $D(p)$ en produit de polynômes irréductibles\footnote{Nous 
rappelons que dans $\mathbb{R}[p]$, les polynômes irréductibles sont 
les polynômes de degré 1 et les polynômes de degré 2 de discriminant négatif.}
unitaire\footnote{Un polynôme unitaire est un polynôme dont le coefficient 
de degré le plus grand est 1.} de la forme:
\[
D(p)=a\prod_{k=1}^r(p-\alpha_k)^{m_k}\prod_{l=1}^s(p^2+\beta_lp+\gamma_l)^{n_l}
\]
où $a$ est une constante qui est le coefficient du terme de plus haut degré 
de $D(p)$, les $\alpha_k$ sont les pôles de multiplicités $m_k$, les polynômes
de degré 2 sont sans pôles réels (c.-à-d. $\beta_l-4\gamma_l<0$).

Il existe alors une famille unique de réels $A_{k,i}$, $B_{l,j}$ et 
$C_{l,j}$ telles que :
%-------------------------------------------------------------------------------
\begin{align}
S(p)=\sum_{k=1}^r\sum_{i=1}^{m_k} \dfrac{A_{k,i}}{\left(p-\alpha_k\right)^i}+
     \sum_{l=1}^s\sum_{j=1}^{n_l} \dfrac{B_{l,j}p+C_{l,j}}{\left(p^2+\beta_lp+
     \gamma_l\right)^j}
\end{align}
%-------------------------------------------------------------------------------
On appelle cette écriture la \textbf{décomposition en éléments simples} de 
$S(p)$ sur $\mathbb{R}$.
\newcounter{exemple}
\newcommand{\exemple}{%
    \refstepcounter{exemple}%
    \paragraph{Exemple \theexemple}%
}
%%%%%%%%%%%%%%%%%%%%%%%%%%%%%%%%%%%%%%%%%%%%%%%%%%%%%%%%%%%%%%%%%%%%%%%%%%%%%%%%
\exemple
%%%%%%%%%%%%%%%%%%%%%%%%%%%%%%%%%%%%%%%%%%%%%%%%%%%%%%%%%%%%%%%%%%%%%%%%%%%%%%%%
Soit $S(p)$ tel que:
\[
S(p)=\dfrac{1}{(p+1)(p+2)}
\]
La décomposition en éléments simples de $S(p)$ est donnée par :
\[
    S(p)=\dfrac{A}{p+1}+\dfrac{B}{p+2}
\]
%%%%%%%%%%%%%%%%%%%%%%%%%%%%%%%%%%%%%%%%%%%%%%%%%%%%%%%%%%%%%%%%%%%%%%%%%%%%%%%%
\exemple
%%%%%%%%%%%%%%%%%%%%%%%%%%%%%%%%%%%%%%%%%%%%%%%%%%%%%%%%%%%%%%%%%%%%%%%%%%%%%%%%
Soit $S(p)$ tel que:
\[
S(p)=\dfrac{1}{(p+1)(p+2)^2}
\]
La décomposition en éléments simples de $S(p)$ est donnée par :
\[
    S(p)=\dfrac{A}{p+1}+\dfrac{B}{p+2}+\dfrac{C}{(p+2)^2}
\]
%%%%%%%%%%%%%%%%%%%%%%%%%%%%%%%%%%%%%%%%%%%%%%%%%%%%%%%%%%%%%%%%%%%%%%%%%%%%%%%%
\exemple
%%%%%%%%%%%%%%%%%%%%%%%%%%%%%%%%%%%%%%%%%%%%%%%%%%%%%%%%%%%%%%%%%%%%%%%%%%%%%%%%
Soit $S(p)$ tel que:
\[
S(p)=\dfrac{1}{(p^2-1)(p^2+1)^2}
\]
où $D(p)$ se factorise sous la forme :
\[
D(p)=(p^2-1)(p^2+1)^2=(p-1)(p+1)(p^2+1)^2
\]
On obtient une décomposition en éléments simples de $S(p)$ de la forme :
\[
    S(p)=\dfrac{A}{p-1}+\dfrac{B}{p+1}+\dfrac{Cp+D}{p^2+1}+\dfrac{Ep+F}{(p^2+1)^2}
\]
%%%%%%%%%%%%%%%%%%%%%%%%%%%%%%%%%%%%%%%%%%%%%%%%%%%%%%%%%%%%%%%%%%%%%%%%%%%%%%%%
\exemple
%%%%%%%%%%%%%%%%%%%%%%%%%%%%%%%%%%%%%%%%%%%%%%%%%%%%%%%%%%%%%%%%%%%%%%%%%%%%%%%%
Soit $S(p)$ tel que:
\[
S(p)=\dfrac{4p^3}{(p^2-1)^2}
\]
où $D(p)$ se factorise sous la forme :
\[
D(p)=(p^2-1)^2=\left((p-1)(p+1)\right)^2
\]
On obtient une décomposition en éléments simples de $S(p)$ de la forme :
\[
S(p)=\dfrac{A}{p-1}+\dfrac{B}{(p-1)^2}+\dfrac{C}{p+1}+\dfrac{D}{(p+1)^2}
\]
%%%%%%%%%%%%%%%%%%%%%%%%%%%%%%%%%%%%%%%%%%%%%%%%%%%%%%%%%%%%%%%%%%%%%%%%%%%%%%%%
%%%%%%%%%%%%%%%%%%%%%%%%%%%%%%%%%%%%%%%%%%%%%%%%%%%%%%%%%%%%%%%%%%%%%%%%%%%%%%%%
%%%%%%%%%%%%%%%%%%%%%%%%%%%%%%%%%%%%%%%%%%%%%%%%%%%%%%%%%%%%%%%%%%%%%%%%%%%%%%%%
\section[Détermination des coefficients de la DES]
        {Détermination des coefficients de la 
         décomposition en éléments simples}
%%%%%%%%%%%%%%%%%%%%%%%%%%%%%%%%%%%%%%%%%%%%%%%%%%%%%%%%%%%%%%%%%%%%%%%%%%%%%%%%
%%%%%%%%%%%%%%%%%%%%%%%%%%%%%%%%%%%%%%%%%%%%%%%%%%%%%%%%%%%%%%%%%%%%%%%%%%%%%%%%
%%%%%%%%%%%%%%%%%%%%%%%%%%%%%%%%%%%%%%%%%%%%%%%%%%%%%%%%%%%%%%%%%%%%%%%%%%%%%%%%
Après avoir obtenu la~\gls{des}, il faut alors en déterminer les coefficients.
Il existe plusieurs méthodes permettant cette détermination, toutes donnant 
évidemment le même résultat. Le choix de la méthode dépend de la décomposition.
On considérera alors la méthode la plus efficace et rapide.

%%%%%%%%%%%%%%%%%%%%%%%%%%%%%%%%%%%%%%%%%%%%%%%%%%%%%%%%%%%%%%%%%%%%%%%%%%%%%%%%
%%%%%%%%%%%%%%%%%%%%%%%%%%%%%%%%%%%%%%%%%%%%%%%%%%%%%%%%%%%%%%%%%%%%%%%%%%%%%%%%
\subsection{Par identification}
%%%%%%%%%%%%%%%%%%%%%%%%%%%%%%%%%%%%%%%%%%%%%%%%%%%%%%%%%%%%%%%%%%%%%%%%%%%%%%%%
%%%%%%%%%%%%%%%%%%%%%%%%%%%%%%%%%%%%%%%%%%%%%%%%%%%%%%%%%%%%%%%%%%%%%%%%%%%%%%%%
Cette méthode est toujours applicable, mais pourra s'avérer être fastidieuse
à mettre en œuvre dans certains cas.
Celle-ci consiste à réduire la décomposition au même dénominateur et à établir
un système d'équations sur les coefficients à déterminer.
\setcounter{exemple}{0}
%%%%%%%%%%%%%%%%%%%%%%%%%%%%%%%%%%%%%%%%%%%%%%%%%%%%%%%%%%%%%%%%%%%%%%%%%%%%%%%%
\exemple
%%%%%%%%%%%%%%%%%%%%%%%%%%%%%%%%%%%%%%%%%%%%%%%%%%%%%%%%%%%%%%%%%%%%%%%%%%%%%%%%
Soit $S(p)$ tel que:
\[
S(p)=\dfrac{1}{(p+1)(p+2)}
\]
sa décomposition en éléments simples est donnée par :
\[
S(p)=\dfrac{A}{p+1}+\dfrac{B}{p+2}
\]
au même dénominateur on obtient
\[
S(p)=\dfrac{A(p+2)+B(p+1)}{(p+1)(p+2)}
    =\dfrac{(A+B)p+(2A+B)}{(p+1)(p+2)}
    =\dfrac{1}{(p+1)(p+2)}
\]
Le système d'équations sur les coefficients de la décomposition est :
\[
\begin{cases}
    A+B=0\\
    2A+B=1
\end{cases}
\]
Sa résolution nous donne, $A=1$ et $B=-1$.
La décomposition en éléments simples de $S(p)$ est donc finalement donnée 
par :
\[
S(p)=\dfrac{1}{p+1}-\dfrac{1}{p+2}
\]
%%%%%%%%%%%%%%%%%%%%%%%%%%%%%%%%%%%%%%%%%%%%%%%%%%%%%%%%%%%%%%%%%%%%%%%%%%%%%%%%
\exemple
%%%%%%%%%%%%%%%%%%%%%%%%%%%%%%%%%%%%%%%%%%%%%%%%%%%%%%%%%%%%%%%%%%%%%%%%%%%%%%%%
Soit $S(p)$ tel que:
\[
S(p)=\dfrac{1}{(p+1)(p+2)^2}
\]
On obtient une décomposition en éléments simples de $S(p)$ de la forme :
\[
    S(p)=\dfrac{A}{p+1}+\dfrac{B}{p+2}+\dfrac{C}{(p+2)^2}
\]
au même dénominateur on obtient :
\[
    S(p)=\dfrac{(A+B)p^2+(4A+3B+C)p+(4A+2B+C)}{(p+1)(p+2)^2}
        =\dfrac{1}{(p+1)(p+2)^2}
\]
Le système d'équations sur les coefficients de la décomposition est :
\[
\begin{cases}
    A+B=0\\
    4A+3B+C=0\\
    4A+2B+C=1
\end{cases}
\]
Sa résolution nous donne, $A=1$ et $B=-1$ et $C=-1$.
La décomposition en éléments simples de $S(p)$ est donc finalement donnée 
par :
\[
    S(p)=\dfrac{1}{p+1}-\dfrac{1}{p+2}-\dfrac{1}{(p+2)^2}
\]

%%%%%%%%%%%%%%%%%%%%%%%%%%%%%%%%%%%%%%%%%%%%%%%%%%%%%%%%%%%%%%%%%%%%%%%%%%%%%%%%
%%%%%%%%%%%%%%%%%%%%%%%%%%%%%%%%%%%%%%%%%%%%%%%%%%%%%%%%%%%%%%%%%%%%%%%%%%%%%%%%
%\subsection{Multiplication/Substitution}
%%%%%%%%%%%%%%%%%%%%%%%%%%%%%%%%%%%%%%%%%%%%%%%%%%%%%%%%%%%%%%%%%%%%%%%%%%%%%%%%
%%%%%%%%%%%%%%%%%%%%%%%%%%%%%%%%%%%%%%%%%%%%%%%%%%%%%%%%%%%%%%%%%%%%%%%%%%%%%%%%

%%%%%%%%%%%%%%%%%%%%%%%%%%%%%%%%%%%%%%%%%%%%%%%%%%%%%%%%%%%%%%%%%%%%%%%%%%%%%%%%
%%%%%%%%%%%%%%%%%%%%%%%%%%%%%%%%%%%%%%%%%%%%%%%%%%%%%%%%%%%%%%%%%%%%%%%%%%%%%%%%
%\subsection{\'Evaluation}
%%%%%%%%%%%%%%%%%%%%%%%%%%%%%%%%%%%%%%%%%%%%%%%%%%%%%%%%%%%%%%%%%%%%%%%%%%%%%%%%
%%%%%%%%%%%%%%%%%%%%%%%%%%%%%%%%%%%%%%%%%%%%%%%%%%%%%%%%%%%%%%%%%%%%%%%%%%%%%%%%
%\acpl
%%%%%%%%%%%%%%%%%%%%%%%%%%%%%%%%%%%%%%%%%%%%%%%%%%%%%%%%%%%%%%%%%%%%%%%%%%%%%%%%
%%%%%%%%%%%%%%%%%%%%%%%%%%%%%%%%%%%%%%%%%%%%%%%%%%%%%%%%%%%%%%%%%%%%%%%%%%%%%%%%
%\subsection{Parité}
%%%%%%%%%%%%%%%%%%%%%%%%%%%%%%%%%%%%%%%%%%%%%%%%%%%%%%%%%%%%%%%%%%%%%%%%%%%%%%%%
%%%%%%%%%%%%%%%%%%%%%%%%%%%%%%%%%%%%%%%%%%%%%%%%%%%%%%%%%%%%%%%%%%%%%%%%%%%%%%%%
%\acpl
%%%%%%%%%%%%%%%%%%%%%%%%%%%%%%%%%%%%%%%%%%%%%%%%%%%%%%%%%%%%%%%%%%%%%%%%%%%%%%%%
%%%%%%%%%%%%%%%%%%%%%%%%%%%%%%%%%%%%%%%%%%%%%%%%%%%%%%%%%%%%%%%%%%%%%%%%%%%%%%%%
%\subsection{Passage à la limite}
%%%%%%%%%%%%%%%%%%%%%%%%%%%%%%%%%%%%%%%%%%%%%%%%%%%%%%%%%%%%%%%%%%%%%%%%%%%%%%%%
%%%%%%%%%%%%%%%%%%%%%%%%%%%%%%%%%%%%%%%%%%%%%%%%%%%%%%%%%%%%%%%%%%%%%%%%%%%%%%%%
%\acpl
%%%%%%%%%%%%%%%%%%%%%%%%%%%%%%%%%%%%%%%%%%%%%%%%%%%%%%%%%%%%%%%%%%%%%%%%%%%%%%%%
%%%%%%%%%%%%%%%%%%%%%%%%%%%%%%%%%%%%%%%%%%%%%%%%%%%%%%%%%%%%%%%%%%%%%%%%%%%%%%%%
%%%%%%%%%%%%%%%%%%%%%%%%%%%%%%%%%%%%%%%%%%%%%%%%%%%%%%%%%%%%%%%%%%%%%%%%%%%%%%%%
%%%%%%%%%%%%%%%%%%%%%%%%%%%%%%%%%%%%%%%%%%%%%%%%%%%%%%%%%%%%%%%%%%%%%%%%%%%%%%%%
%annexe_des.tex
